% Options for packages loaded elsewhere
% Options for packages loaded elsewhere
\PassOptionsToPackage{unicode}{hyperref}
\PassOptionsToPackage{hyphens}{url}
\PassOptionsToPackage{dvipsnames,svgnames,x11names}{xcolor}
%
\documentclass[
  spanish,
  letterpaper,
  DIV=11,
  numbers=noendperiod]{scrartcl}
\usepackage{xcolor}
\usepackage{amsmath,amssymb}
\setcounter{secnumdepth}{-\maxdimen} % remove section numbering
\usepackage{iftex}
\ifPDFTeX
  \usepackage[T1]{fontenc}
  \usepackage[utf8]{inputenc}
  \usepackage{textcomp} % provide euro and other symbols
\else % if luatex or xetex
  \usepackage{unicode-math} % this also loads fontspec
  \defaultfontfeatures{Scale=MatchLowercase}
  \defaultfontfeatures[\rmfamily]{Ligatures=TeX,Scale=1}
\fi
\usepackage{lmodern}
\ifPDFTeX\else
  % xetex/luatex font selection
\fi
% Use upquote if available, for straight quotes in verbatim environments
\IfFileExists{upquote.sty}{\usepackage{upquote}}{}
\IfFileExists{microtype.sty}{% use microtype if available
  \usepackage[]{microtype}
  \UseMicrotypeSet[protrusion]{basicmath} % disable protrusion for tt fonts
}{}
\makeatletter
\@ifundefined{KOMAClassName}{% if non-KOMA class
  \IfFileExists{parskip.sty}{%
    \usepackage{parskip}
  }{% else
    \setlength{\parindent}{0pt}
    \setlength{\parskip}{6pt plus 2pt minus 1pt}}
}{% if KOMA class
  \KOMAoptions{parskip=half}}
\makeatother
% Make \paragraph and \subparagraph free-standing
\makeatletter
\ifx\paragraph\undefined\else
  \let\oldparagraph\paragraph
  \renewcommand{\paragraph}{
    \@ifstar
      \xxxParagraphStar
      \xxxParagraphNoStar
  }
  \newcommand{\xxxParagraphStar}[1]{\oldparagraph*{#1}\mbox{}}
  \newcommand{\xxxParagraphNoStar}[1]{\oldparagraph{#1}\mbox{}}
\fi
\ifx\subparagraph\undefined\else
  \let\oldsubparagraph\subparagraph
  \renewcommand{\subparagraph}{
    \@ifstar
      \xxxSubParagraphStar
      \xxxSubParagraphNoStar
  }
  \newcommand{\xxxSubParagraphStar}[1]{\oldsubparagraph*{#1}\mbox{}}
  \newcommand{\xxxSubParagraphNoStar}[1]{\oldsubparagraph{#1}\mbox{}}
\fi
\makeatother

\usepackage{color}
\usepackage{fancyvrb}
\newcommand{\VerbBar}{|}
\newcommand{\VERB}{\Verb[commandchars=\\\{\}]}
\DefineVerbatimEnvironment{Highlighting}{Verbatim}{commandchars=\\\{\}}
% Add ',fontsize=\small' for more characters per line
\usepackage{framed}
\definecolor{shadecolor}{RGB}{241,243,245}
\newenvironment{Shaded}{\begin{snugshade}}{\end{snugshade}}
\newcommand{\AlertTok}[1]{\textcolor[rgb]{0.68,0.00,0.00}{#1}}
\newcommand{\AnnotationTok}[1]{\textcolor[rgb]{0.37,0.37,0.37}{#1}}
\newcommand{\AttributeTok}[1]{\textcolor[rgb]{0.40,0.45,0.13}{#1}}
\newcommand{\BaseNTok}[1]{\textcolor[rgb]{0.68,0.00,0.00}{#1}}
\newcommand{\BuiltInTok}[1]{\textcolor[rgb]{0.00,0.23,0.31}{#1}}
\newcommand{\CharTok}[1]{\textcolor[rgb]{0.13,0.47,0.30}{#1}}
\newcommand{\CommentTok}[1]{\textcolor[rgb]{0.37,0.37,0.37}{#1}}
\newcommand{\CommentVarTok}[1]{\textcolor[rgb]{0.37,0.37,0.37}{\textit{#1}}}
\newcommand{\ConstantTok}[1]{\textcolor[rgb]{0.56,0.35,0.01}{#1}}
\newcommand{\ControlFlowTok}[1]{\textcolor[rgb]{0.00,0.23,0.31}{\textbf{#1}}}
\newcommand{\DataTypeTok}[1]{\textcolor[rgb]{0.68,0.00,0.00}{#1}}
\newcommand{\DecValTok}[1]{\textcolor[rgb]{0.68,0.00,0.00}{#1}}
\newcommand{\DocumentationTok}[1]{\textcolor[rgb]{0.37,0.37,0.37}{\textit{#1}}}
\newcommand{\ErrorTok}[1]{\textcolor[rgb]{0.68,0.00,0.00}{#1}}
\newcommand{\ExtensionTok}[1]{\textcolor[rgb]{0.00,0.23,0.31}{#1}}
\newcommand{\FloatTok}[1]{\textcolor[rgb]{0.68,0.00,0.00}{#1}}
\newcommand{\FunctionTok}[1]{\textcolor[rgb]{0.28,0.35,0.67}{#1}}
\newcommand{\ImportTok}[1]{\textcolor[rgb]{0.00,0.46,0.62}{#1}}
\newcommand{\InformationTok}[1]{\textcolor[rgb]{0.37,0.37,0.37}{#1}}
\newcommand{\KeywordTok}[1]{\textcolor[rgb]{0.00,0.23,0.31}{\textbf{#1}}}
\newcommand{\NormalTok}[1]{\textcolor[rgb]{0.00,0.23,0.31}{#1}}
\newcommand{\OperatorTok}[1]{\textcolor[rgb]{0.37,0.37,0.37}{#1}}
\newcommand{\OtherTok}[1]{\textcolor[rgb]{0.00,0.23,0.31}{#1}}
\newcommand{\PreprocessorTok}[1]{\textcolor[rgb]{0.68,0.00,0.00}{#1}}
\newcommand{\RegionMarkerTok}[1]{\textcolor[rgb]{0.00,0.23,0.31}{#1}}
\newcommand{\SpecialCharTok}[1]{\textcolor[rgb]{0.37,0.37,0.37}{#1}}
\newcommand{\SpecialStringTok}[1]{\textcolor[rgb]{0.13,0.47,0.30}{#1}}
\newcommand{\StringTok}[1]{\textcolor[rgb]{0.13,0.47,0.30}{#1}}
\newcommand{\VariableTok}[1]{\textcolor[rgb]{0.07,0.07,0.07}{#1}}
\newcommand{\VerbatimStringTok}[1]{\textcolor[rgb]{0.13,0.47,0.30}{#1}}
\newcommand{\WarningTok}[1]{\textcolor[rgb]{0.37,0.37,0.37}{\textit{#1}}}

\usepackage{longtable,booktabs,array}
\newcounter{none} % for unnumbered tables
\usepackage{calc} % for calculating minipage widths
% Correct order of tables after \paragraph or \subparagraph
\usepackage{etoolbox}
\makeatletter
\patchcmd\longtable{\par}{\if@noskipsec\mbox{}\fi\par}{}{}
\makeatother
% Allow footnotes in longtable head/foot
\IfFileExists{footnotehyper.sty}{\usepackage{footnotehyper}}{\usepackage{footnote}}
\makesavenoteenv{longtable}
\usepackage{graphicx}
\makeatletter
\newsavebox\pandoc@box
\newcommand*\pandocbounded[1]{% scales image to fit in text height/width
  \sbox\pandoc@box{#1}%
  \Gscale@div\@tempa{\textheight}{\dimexpr\ht\pandoc@box+\dp\pandoc@box\relax}%
  \Gscale@div\@tempb{\linewidth}{\wd\pandoc@box}%
  \ifdim\@tempb\p@<\@tempa\p@\let\@tempa\@tempb\fi% select the smaller of both
  \ifdim\@tempa\p@<\p@\scalebox{\@tempa}{\usebox\pandoc@box}%
  \else\usebox{\pandoc@box}%
  \fi%
}
% Set default figure placement to htbp
\def\fps@figure{htbp}
\makeatother

\ifLuaTeX
  \usepackage{luacolor}
  \usepackage[soul]{lua-ul}
\else
  \usepackage{soul}
\fi


\ifLuaTeX
\usepackage[bidi=basic,shorthands=off]{babel}
\else
\usepackage[bidi=default,shorthands=off]{babel}
\fi
\ifLuaTeX
  \usepackage{selnolig} % disable illegal ligatures
\fi


\setlength{\emergencystretch}{3em} % prevent overfull lines

\providecommand{\tightlist}{%
  \setlength{\itemsep}{0pt}\setlength{\parskip}{0pt}}



 


<style>
body {
  text-align: justify;
}
</style>
\KOMAoption{captions}{tableheading}
\makeatletter
\@ifpackageloaded{caption}{}{\usepackage{caption}}
\AtBeginDocument{%
\ifdefined\contentsname
  \renewcommand*\contentsname{Tabla de contenidos}
\else
  \newcommand\contentsname{Tabla de contenidos}
\fi
\ifdefined\listfigurename
  \renewcommand*\listfigurename{Listado de Figuras}
\else
  \newcommand\listfigurename{Listado de Figuras}
\fi
\ifdefined\listtablename
  \renewcommand*\listtablename{Listado de Tablas}
\else
  \newcommand\listtablename{Listado de Tablas}
\fi
\ifdefined\figurename
  \renewcommand*\figurename{Figura}
\else
  \newcommand\figurename{Figura}
\fi
\ifdefined\tablename
  \renewcommand*\tablename{Tabla}
\else
  \newcommand\tablename{Tabla}
\fi
}
\@ifpackageloaded{float}{}{\usepackage{float}}
\floatstyle{ruled}
\@ifundefined{c@chapter}{\newfloat{codelisting}{h}{lop}}{\newfloat{codelisting}{h}{lop}[chapter]}
\floatname{codelisting}{Listado}
\newcommand*\listoflistings{\listof{codelisting}{Listado de Listados}}
\makeatother
\makeatletter
\makeatother
\makeatletter
\@ifpackageloaded{caption}{}{\usepackage{caption}}
\@ifpackageloaded{subcaption}{}{\usepackage{subcaption}}
\makeatother
\usepackage{bookmark}
\IfFileExists{xurl.sty}{\usepackage{xurl}}{} % add URL line breaks if available
\urlstyle{same}
\hypersetup{
  pdftitle={Data Wrangling en R 2},
  pdflang={es},
  colorlinks=true,
  linkcolor={blue},
  filecolor={Maroon},
  citecolor={Blue},
  urlcolor={Blue},
  pdfcreator={LaTeX via pandoc}}


\title{Data Wrangling en R 2}
\author{}
\date{2026-04-07}
\begin{document}
\maketitle

\renewcommand*\contentsname{Tabla de contenidos}
{
\setcounter{tocdepth}{3}
\tableofcontents
}

\subsection{De qué trata la clase de
hoy}\label{de-quuxe9-trata-la-clase-de-hoy}

Hoy vamos a continuar con los temas de la clase pasada, pero vamos a
utilizar otras herramientas que nos van a facilitar muchas de las tareas
que veníamos haciendo.

\begin{enumerate}
\def\labelenumi{\arabic{enumi}.}
\item
  \textbf{Primera parte}: vamos a ver las consignas del TP.
\item
  \textbf{Segunda parte}: vamos a revisar el final de la clase anterior,
  que quedó inconcluso y tenía algunos bugs.
\item
  \textbf{Tercera parte}: vamos a ver otra base de datos. La idea en
  este caso es volver a hacer algunas de las operaciones que hicimos
  antes valiéndonos de \texttt{janitor} y \texttt{stringi}, que son
  librerías que sintetizan operaciones de limpieza que resultan
  engorrosas de escribir a mano. La idea es que este \emph{notebook y}
  el anterior les sirva de machete para que puedan hacer más
  tranquilos/as el TP individual. Vamos a ver, además, cómo ordenar
  variables en forma jerárquica.
\item
  \textbf{Cuarta parte}: seguramente la clase que viene veremos cómo
  organizar \emph{scripts} dentro de un proyecto grande para que cada
  archivo realice una tarea en particular.
\end{enumerate}

\subsection{3. Tercera Parte}\label{tercera-parte}

Vamos a trabajar con datos del \textbf{SINTA (Sistema de Inteligencia
Turística Argentina)}, que es el portal de datos abiertos de la
Subsecretaría de Turismo. Utilizaremos específicamente dos bases de
datos de la \textbf{Encuesta de Turismo Internacional (ETI)}: una
correspondiente al
\href{https://datos.yvera.gob.ar/dataset/encuesta-turismo-internacional/archivo/dae8b7f8-f3dc-431b-887e-3de4b54f09a9}{\textbf{Aeropuerto
Internacional de Córdoba}} y otra que agrupa el movimiento de
\href{https://datos.yvera.gob.ar/dataset/encuesta-turismo-internacional/archivo/32cd65c4-7558-48cf-a8ac-bbb07147b5a1}{\textbf{CABA}
\textbf{y PBA} a través de los arribos en \textbf{Ezeiza y Aeroparque}}.

Estas bases nos permiten analizar el volumen de \textbf{turistas no
residentes}, sus \textbf{pernoctes} y la \textbf{estadía media} según su
país de residencia, lo que nos servirá de insumo para practicar técnicas
de limpieza, normalización y transformación de datos.

\ul{En este caso no vamos a trabajar sobre una hipótesis particular
porque armé el análisis para que se puedan realizar las operaciones que
nos faltan revisar.}

\subsubsection{Librerías y los datos}\label{libreruxedas-y-los-datos}

\begin{Shaded}
\begin{Highlighting}[]
\CommentTok{\# install.packages(c("tidyverse", "janitor"))}
\FunctionTok{library}\NormalTok{(tidyverse)}
\end{Highlighting}
\end{Shaded}

\begin{verbatim}
Warning: package 'tidyverse' was built under R version 4.5.3
\end{verbatim}

\begin{verbatim}
-- Attaching core tidyverse packages ------------------------ tidyverse 2.0.0 --
v dplyr     1.1.4     v readr     2.1.6
v forcats   1.0.1     v stringr   1.6.0
v ggplot2   4.0.1     v tibble    3.3.0
v lubridate 1.9.4     v tidyr     1.3.2
v purrr     1.2.0     
-- Conflicts ------------------------------------------ tidyverse_conflicts() --
x dplyr::filter() masks stats::filter()
x dplyr::lag()    masks stats::lag()
i Use the conflicted package (<http://conflicted.r-lib.org/>) to force all conflicts to become errors
\end{verbatim}

\begin{Shaded}
\begin{Highlighting}[]
\FunctionTok{library}\NormalTok{(janitor)}
\end{Highlighting}
\end{Shaded}

\begin{verbatim}
Warning: package 'janitor' was built under R version 4.5.3
\end{verbatim}

\begin{verbatim}

Attaching package: 'janitor'

The following objects are masked from 'package:stats':

    chisq.test, fisher.test
\end{verbatim}

\begin{Shaded}
\begin{Highlighting}[]
\FunctionTok{library}\NormalTok{(stringi)}
\end{Highlighting}
\end{Shaded}

\begin{Shaded}
\begin{Highlighting}[]
\NormalTok{path\_ba }\OtherTok{=} \StringTok{"../data/estadias\_pba\_caba.csv"}
\NormalTok{estadias\_ba }\OtherTok{=} \FunctionTok{read\_csv}\NormalTok{(}
\NormalTok{  path\_ba, }
  \CommentTok{\# skip=0 si tiene comentarios en las primeras antes de la tabla, este parámetro  indica la cantidad de lineas que ocupan, para saltarlas}
\NormalTok{  )}
\end{Highlighting}
\end{Shaded}

\begin{verbatim}
Rows: 432 Columns: 6
-- Column specification --------------------------------------------------------
Delimiter: ","
chr  (2): país de residencia, observaciones
dbl  (3): turistas_no_residentes, pernoctes_no_residentes, estadia_media_no_...
date (1): indice_tiempo

i Use `spec()` to retrieve the full column specification for this data.
i Specify the column types or set `show_col_types = FALSE` to quiet this message.
\end{verbatim}

\begin{Shaded}
\begin{Highlighting}[]
\FunctionTok{head}\NormalTok{(estadias\_ba)}
\end{Highlighting}
\end{Shaded}

\begin{verbatim}
# A tibble: 6 x 6
  indice_tiempo `país de residencia` turistas_no_residentes
  <date>        <chr>                                 <dbl>
1 2014-01-01    Bolivia                               17181
2 2014-01-01    Brasil                               149582
3 2014-01-01    Chile                                 48953
4 2014-01-01    Paraguay                              17316
5 2014-01-01    Uruguay                               11429
6 2014-01-01    EE.UU y Canadá                        82090
# i 3 more variables: pernoctes_no_residentes <dbl>,
#   estadia_media_no_residentes <dbl>, observaciones <chr>
\end{verbatim}

\begin{Shaded}
\begin{Highlighting}[]
\NormalTok{path\_co }\OtherTok{=} \StringTok{"../data/estadias\_cordoba.csv"}
\NormalTok{estadias\_co }\OtherTok{=} \FunctionTok{read\_csv}\NormalTok{(path\_co)}
\end{Highlighting}
\end{Shaded}

\begin{verbatim}
Rows: 240 Columns: 6
-- Column specification --------------------------------------------------------
Delimiter: ","
chr  (2): país de residencia, observaciones
dbl  (3): turistas_no_residentes, pernoctes_no_residentes, estadia_media_no_...
date (1): indice_tiempo

i Use `spec()` to retrieve the full column specification for this data.
i Specify the column types or set `show_col_types = FALSE` to quiet this message.
\end{verbatim}

\begin{Shaded}
\begin{Highlighting}[]
\FunctionTok{head}\NormalTok{(estadias\_co)}
\end{Highlighting}
\end{Shaded}

\begin{verbatim}
# A tibble: 6 x 6
  indice_tiempo `país de residencia`     turistas_no_residentes
  <date>        <chr>                                     <dbl>
1 2014-01-01    Brasil                                     1964
2 2014-01-01    Chile                                      3138
3 2014-01-01    EE.UU, Canadá y México                     3635
4 2014-01-01    Resto de América                           2290
5 2014-01-01    Europa y resto del mundo                   2159
6 2014-04-01    Brasil                                     2049
# i 3 more variables: pernoctes_no_residentes <dbl>,
#   estadia_media_no_residentes <dbl>, observaciones <chr>
\end{verbatim}

\subsubsection{Normalización de variables y remoción de datos
faltantes}\label{normalizaciuxf3n-de-variables-y-remociuxf3n-de-datos-faltantes}

Armamos una función que, dado un dataframe, realice ambas operaciones

\begin{Shaded}
\begin{Highlighting}[]
\NormalTok{normalize\_dataframe }\OtherTok{=} \ControlFlowTok{function}\NormalTok{(dataframe, }\AttributeTok{clean =} \ConstantTok{TRUE}\NormalTok{, }\AttributeTok{which\_remove =} \StringTok{"cols"}\NormalTok{) \{}
  \CommentTok{\# Normalizacion}
  \ControlFlowTok{if}\NormalTok{ (clean) \{}
\NormalTok{    data\_clean }\OtherTok{=} \FunctionTok{clean\_names}\NormalTok{(dataframe)}
\NormalTok{  \} }\ControlFlowTok{else}\NormalTok{ \{}
\NormalTok{    data\_clean }\OtherTok{=}\NormalTok{ dataframe}
\NormalTok{  \}}
  
  \CommentTok{\# Eliminacion de datos faltantes}
  \ControlFlowTok{if}\NormalTok{ (which\_remove }\SpecialCharTok{==} \StringTok{"both"}\NormalTok{) \{}
\NormalTok{    data\_clean }\OtherTok{=}\NormalTok{ data\_clean }\SpecialCharTok{|\textgreater{}} 
\NormalTok{      janitor}\SpecialCharTok{::}\FunctionTok{remove\_empty}\NormalTok{(}\AttributeTok{which =} \StringTok{\textquotesingle{}cols\textquotesingle{}}\NormalTok{) }\SpecialCharTok{|\textgreater{}} 
\NormalTok{      janitor}\SpecialCharTok{::}\FunctionTok{remove\_empty}\NormalTok{(}\AttributeTok{which =} \StringTok{\textquotesingle{}rows\textquotesingle{}}\NormalTok{)}
\NormalTok{  \} }\ControlFlowTok{else}\NormalTok{ \{}
\NormalTok{    data\_clean }\OtherTok{=}\NormalTok{ janitor}\SpecialCharTok{::}\FunctionTok{remove\_empty}\NormalTok{(data\_clean, }\AttributeTok{which =}\NormalTok{ which\_remove)}
\NormalTok{  \}}
  
  \FunctionTok{return}\NormalTok{(data\_clean)}
\NormalTok{\}}
\CommentTok{\# normalize\_dataframe(estadias\_co, clean = FALSE, which\_remove = \textquotesingle{}both\textquotesingle{})}
\end{Highlighting}
\end{Shaded}

Esta función recibe como input un dataframe, normaliza los nombres de
las variables (columnas) y, luego, remueve las observaciones
(\texttt{rows}) o variables (\texttt{cols}) que tengan \ul{todos} sus
valores faltantes (\texttt{NA}).

Para llamar una función que está definida en dos librerías diferentes
con el mismo nombre, podemos escribir primero la librería creadora para
desambiguar el conflicto ---\texttt{janitor::clean\_names()}. En este
ejemplo especificar la librería padre está de más porque
\texttt{clean\_names} no se define en otras librerías.

En caso de querer desambiguar una función \emph{built-in} (pre definida
en R), podemos hacerlo así: \texttt{base::una\_funcion()}.

\begin{Shaded}
\begin{Highlighting}[]
\NormalTok{estadias\_co\_clean }\OtherTok{=} \FunctionTok{normalize\_dataframe}\NormalTok{(estadias\_co, }\AttributeTok{which\_remove =} \StringTok{\textquotesingle{}both\textquotesingle{}}\NormalTok{)}
\NormalTok{estadias\_ba\_clean }\OtherTok{=} \FunctionTok{normalize\_dataframe}\NormalTok{(estadias\_ba, }\AttributeTok{which\_remove =} \StringTok{"both"}\NormalTok{)}
\FunctionTok{cat}\NormalTok{(}
  \FunctionTok{c}\NormalTok{(}\StringTok{\textquotesingle{}Estadías Córdoba antes:\textquotesingle{}}\NormalTok{, }\FunctionTok{nrow}\NormalTok{(estadias\_co), }\StringTok{\textquotesingle{}y después:\textquotesingle{}}\NormalTok{, }\FunctionTok{nrow}\NormalTok{(estadias\_co\_clean))}
\NormalTok{)}
\end{Highlighting}
\end{Shaded}

\begin{verbatim}
Estadías Córdoba antes: 240 y después: 239
\end{verbatim}

\begin{Shaded}
\begin{Highlighting}[]
\FunctionTok{cat}\NormalTok{(}
  \FunctionTok{c}\NormalTok{(}\StringTok{\textquotesingle{}}\SpecialCharTok{\textbackslash{}n}\StringTok{Estadías CABA antes:\textquotesingle{}}\NormalTok{, }\FunctionTok{nrow}\NormalTok{(estadias\_ba), }\StringTok{\textquotesingle{}y después:\textquotesingle{}}\NormalTok{, }\FunctionTok{nrow}\NormalTok{(estadias\_ba\_clean))}
\NormalTok{)}
\end{Highlighting}
\end{Shaded}

\begin{verbatim}

Estadías CABA antes: 432 y después: 430
\end{verbatim}

Notemos que hay países que no comparten las provincias, entonces es
conveniente hacer un \texttt{inner\_join}.

\begin{Shaded}
\begin{Highlighting}[]
\NormalTok{estadias\_ba\_clean }\SpecialCharTok{|\textgreater{}} \FunctionTok{select}\NormalTok{(pais\_de\_residencia) }\SpecialCharTok{|\textgreater{}} \FunctionTok{unique}\NormalTok{()}
\end{Highlighting}
\end{Shaded}

\begin{verbatim}
# A tibble: 9 x 1
  pais_de_residencia
  <chr>             
1 Bolivia           
2 Brasil            
3 Chile             
4 Paraguay          
5 Uruguay           
6 EE.UU y Canadá    
7 Resto América     
8 Europa            
9 Resto del Mundo   
\end{verbatim}

\begin{Shaded}
\begin{Highlighting}[]
\NormalTok{estadias\_co\_clean }\SpecialCharTok{|\textgreater{}} \FunctionTok{select}\NormalTok{(pais\_de\_residencia) }\SpecialCharTok{|\textgreater{}} \FunctionTok{unique}\NormalTok{()}
\end{Highlighting}
\end{Shaded}

\begin{verbatim}
# A tibble: 5 x 1
  pais_de_residencia      
  <chr>                   
1 Brasil                  
2 Chile                   
3 EE.UU, Canadá y México  
4 Resto de América        
5 Europa y resto del mundo
\end{verbatim}

Pero hay que tener cuidado ya que hay países que significan lo mismo,
pero tienen nombres diferentes. Por ejemplo, ``Resto de América'' y
``Resto América''. En este caso, podemos revisar las categorías a mano
dado que son pocas, de todos modos hagamos un tratamiento general.

\begin{Shaded}
\begin{Highlighting}[]
\NormalTok{estadias\_ba\_clean }\OtherTok{=}\NormalTok{ estadias\_ba\_clean }\SpecialCharTok{|\textgreater{}} 
  \FunctionTok{mutate}\NormalTok{(}
    \AttributeTok{pais\_de\_residencia =}\NormalTok{ pais\_de\_residencia }\SpecialCharTok{|\textgreater{}} 
           \FunctionTok{str\_to\_lower}\NormalTok{() }\SpecialCharTok{|\textgreater{}}          \CommentTok{\# todo a minúsculas}
           \FunctionTok{stri\_trans\_general}\NormalTok{(}\StringTok{"Latin{-}ASCII"}\NormalTok{) }\SpecialCharTok{|\textgreater{}} \CommentTok{\# quita tildes }
           \FunctionTok{str\_replace\_all}\NormalTok{(}\StringTok{" "}\NormalTok{, }\StringTok{"\_"}\NormalTok{) }\SpecialCharTok{|\textgreater{}}     \CommentTok{\# espacios por guiones}
           \FunctionTok{str\_remove\_all}\NormalTok{(}\StringTok{"[[:punct:]]"}\NormalTok{) }\SpecialCharTok{|\textgreater{}} \CommentTok{\# quita puntos y signos restantes}
           \FunctionTok{str\_trim}\NormalTok{() }\CommentTok{\# elimina espacios no detectados}
\NormalTok{  ) }\SpecialCharTok{|\textgreater{}} 
  \FunctionTok{mutate}\NormalTok{(}
    \AttributeTok{pais\_de\_residencia =} \FunctionTok{case\_when}\NormalTok{(}
      \FunctionTok{str\_starts}\NormalTok{(pais\_de\_residencia, }\StringTok{"eeuu"}\NormalTok{) }\SpecialCharTok{|} 
      \FunctionTok{str\_starts}\NormalTok{(pais\_de\_residencia, }\StringTok{"euro"}\NormalTok{) }\SpecialCharTok{|} 
      \FunctionTok{str\_starts}\NormalTok{(pais\_de\_residencia, }\StringTok{"resto"}\NormalTok{) }\SpecialCharTok{\textasciitilde{}} \StringTok{"nolimitrofe"}\NormalTok{, }\CommentTok{\# la tilde se usa para separar la condición del resultado}
      \AttributeTok{.default =}\NormalTok{ pais\_de\_residencia }\CommentTok{\# mantiene el nombre original si no es \textquotesingle{}otros\textquotesingle{}}
\NormalTok{    )}
\NormalTok{  )}
\NormalTok{estadias\_co\_clean }\OtherTok{=}\NormalTok{ estadias\_co\_clean }\SpecialCharTok{|\textgreater{}} 
  \FunctionTok{mutate}\NormalTok{(}
    \AttributeTok{pais\_de\_residencia =}\NormalTok{ pais\_de\_residencia }\SpecialCharTok{|\textgreater{}} 
           \FunctionTok{str\_to\_lower}\NormalTok{() }\SpecialCharTok{|\textgreater{}}          \CommentTok{\# todo a minúsculas}
           \FunctionTok{stri\_trans\_general}\NormalTok{(}\StringTok{"Latin{-}ASCII"}\NormalTok{) }\SpecialCharTok{|\textgreater{}} \CommentTok{\# quita tildes }
           \FunctionTok{str\_replace\_all}\NormalTok{(}\StringTok{" "}\NormalTok{, }\StringTok{"\_"}\NormalTok{) }\SpecialCharTok{|\textgreater{}}     \CommentTok{\# espacios por guiones}
           \FunctionTok{str\_remove\_all}\NormalTok{(}\StringTok{"[[:punct:]]"}\NormalTok{) }\SpecialCharTok{|\textgreater{}} \CommentTok{\# quita puntos y signos restantes}
           \FunctionTok{str\_trim}\NormalTok{() }\CommentTok{\# elimina espacios no detectados}
\NormalTok{  ) }\SpecialCharTok{|\textgreater{}} 
  \FunctionTok{mutate}\NormalTok{(}
    \AttributeTok{pais\_de\_residencia =} \FunctionTok{case\_when}\NormalTok{(}
      \FunctionTok{str\_starts}\NormalTok{(pais\_de\_residencia, }\StringTok{"eeuu"}\NormalTok{) }\SpecialCharTok{|} 
      \FunctionTok{str\_starts}\NormalTok{(pais\_de\_residencia, }\StringTok{"euro"}\NormalTok{) }\SpecialCharTok{|} 
      \FunctionTok{str\_starts}\NormalTok{(pais\_de\_residencia, }\StringTok{"resto"}\NormalTok{) }\SpecialCharTok{\textasciitilde{}} \StringTok{"nolimitrofe"}\NormalTok{,}
      \AttributeTok{.default =}\NormalTok{ pais\_de\_residencia }\CommentTok{\# mantiene el nombre original si no es \textquotesingle{}otros\textquotesingle{}}
\NormalTok{    )}
\NormalTok{  )}
\NormalTok{estadias\_ba\_clean }\SpecialCharTok{|\textgreater{}} \FunctionTok{select}\NormalTok{(pais\_de\_residencia) }\SpecialCharTok{|\textgreater{}} \FunctionTok{unique}\NormalTok{()}
\end{Highlighting}
\end{Shaded}

\begin{verbatim}
# A tibble: 6 x 1
  pais_de_residencia
  <chr>             
1 bolivia           
2 brasil            
3 chile             
4 paraguay          
5 uruguay           
6 nolimitrofe       
\end{verbatim}

\begin{Shaded}
\begin{Highlighting}[]
\NormalTok{estadias\_co\_clean }\SpecialCharTok{|\textgreater{}} \FunctionTok{select}\NormalTok{(pais\_de\_residencia) }\SpecialCharTok{|\textgreater{}} \FunctionTok{unique}\NormalTok{()}
\end{Highlighting}
\end{Shaded}

\begin{verbatim}
# A tibble: 3 x 1
  pais_de_residencia
  <chr>             
1 brasil            
2 chile             
3 nolimitrofe       
\end{verbatim}

Ahora sí, notamos que Córdoba registra ingresos por avión de brasileros,
chilenos y países no limítrofes, mientras que BA tiene ingresos más
diversos.

\ul{\textbf{Atención}:} esto no significa que Córdoba no tenga ingresos
de otros países, puede que vuelos no directos no se registren de la
misma forma (y por eso, no aparezcan en la tabla), o que sea conveniente
aterrizar en BA e ir en otro medio de transporte (micro por ej.)

Otro detalle a tener en cuenta es que como agrupamos categorías en
\texttt{pais\_de\_residencia} ahora habrá indices de tiempo repetidos
para un mismo país, lo cual puede conllevar a errores más adelante.

EEUU, Mexico 2022/01/01 --\textgreater{} nolimitrofe 2022/01/01

Mexico 2022/01/01 --\textgreater{} nolimitrofe 2022/01/01

\begin{Shaded}
\begin{Highlighting}[]
\NormalTok{estadias\_ba\_clean }\OtherTok{=}\NormalTok{ estadias\_ba\_clean }\SpecialCharTok{|\textgreater{}} 
  \FunctionTok{group\_by}\NormalTok{(pais\_de\_residencia, indice\_tiempo) }\SpecialCharTok{|\textgreater{}} \CommentTok{\# agrupamos}
    \FunctionTok{summarise}\NormalTok{( }\CommentTok{\# tomamos valores medios /sumas según el indicador}
      \AttributeTok{turistas\_no\_residentes =} \FunctionTok{sum}\NormalTok{(turistas\_no\_residentes, }\AttributeTok{na.rm =} \ConstantTok{TRUE}\NormalTok{),}
      \AttributeTok{pernoctes\_no\_residentes =} \FunctionTok{sum}\NormalTok{(pernoctes\_no\_residentes, }\AttributeTok{na.rm =} \ConstantTok{TRUE}\NormalTok{),}
      \AttributeTok{estadia\_media\_no\_residentes =} \FunctionTok{mean}\NormalTok{(estadia\_media\_no\_residentes, }\AttributeTok{na.rm =} \ConstantTok{TRUE}\NormalTok{),}
      \AttributeTok{.groups =} \StringTok{"drop"} \CommentTok{\# Para que el DF deje de estar agrupado tras el cálculo}
\NormalTok{    )}
\NormalTok{estadias\_co\_clean }\OtherTok{=}\NormalTok{ estadias\_co\_clean }\SpecialCharTok{|\textgreater{}} 
  \FunctionTok{group\_by}\NormalTok{(pais\_de\_residencia, indice\_tiempo) }\SpecialCharTok{|\textgreater{}} \CommentTok{\# agrupamos}
    \FunctionTok{summarise}\NormalTok{( }\CommentTok{\# tomamos valores medios /sumas según el indicador}
      \AttributeTok{turistas\_no\_residentes =} \FunctionTok{sum}\NormalTok{(turistas\_no\_residentes, }\AttributeTok{na.rm =} \ConstantTok{TRUE}\NormalTok{),}
      \AttributeTok{pernoctes\_no\_residentes =} \FunctionTok{sum}\NormalTok{(pernoctes\_no\_residentes, }\AttributeTok{na.rm =} \ConstantTok{TRUE}\NormalTok{),}
      \AttributeTok{estadia\_media\_no\_residentes =} \FunctionTok{mean}\NormalTok{(estadia\_media\_no\_residentes, }\AttributeTok{na.rm =} \ConstantTok{TRUE}\NormalTok{),}
      \AttributeTok{.groups =} \StringTok{"drop"} \CommentTok{\# Para que el DF deje de estar agrupado tras el cálculo}
\NormalTok{    )}
\NormalTok{estadias\_co\_clean }\SpecialCharTok{|\textgreater{}} \FunctionTok{filter}\NormalTok{(pais\_de\_residencia}\SpecialCharTok{==}\StringTok{\textquotesingle{}nolimitrofe\textquotesingle{}}\NormalTok{)}
\end{Highlighting}
\end{Shaded}

\begin{verbatim}
# A tibble: 48 x 5
   pais_de_residencia indice_tiempo turistas_no_residentes
   <chr>              <date>                         <dbl>
 1 nolimitrofe        2014-01-01                      8084
 2 nolimitrofe        2014-04-01                      4163
 3 nolimitrofe        2014-07-01                      7086
 4 nolimitrofe        2014-10-01                      6325
 5 nolimitrofe        2015-01-01                      6239
 6 nolimitrofe        2015-04-01                      5812
 7 nolimitrofe        2015-07-01                      6798
 8 nolimitrofe        2015-10-01                      7808
 9 nolimitrofe        2016-01-01                      8374
10 nolimitrofe        2016-04-01                      5766
# i 38 more rows
# i 2 more variables: pernoctes_no_residentes <dbl>,
#   estadia_media_no_residentes <dbl>
\end{verbatim}

Ahora sí, teniendo todo ordenado podemos unir las dos tablas y ver, como
anticipábamos, que sólo tienen en común Brasil, Chile y no limítrofes.

\begin{Shaded}
\begin{Highlighting}[]
\NormalTok{estadias\_totales }\OtherTok{=}\NormalTok{ estadias\_ba\_clean }\SpecialCharTok{|\textgreater{}} 
  \FunctionTok{inner\_join}\NormalTok{(}
\NormalTok{    estadias\_co\_clean,}
    \AttributeTok{by=}\FunctionTok{c}\NormalTok{(}\StringTok{"pais\_de\_residencia"}\NormalTok{,}\StringTok{"indice\_tiempo"}\NormalTok{),}
    \AttributeTok{suffix =} \FunctionTok{c}\NormalTok{(}\StringTok{"\_BA"}\NormalTok{, }\StringTok{"\_CO"}\NormalTok{) }\CommentTok{\# el sufijo asignado a las tablas izq y derecha (reemplaza el valor defecto que es .x y .y)}
\NormalTok{  ) }\SpecialCharTok{|\textgreater{}} 
  \FunctionTok{select}\NormalTok{(}\SpecialCharTok{!}\FunctionTok{starts\_with}\NormalTok{(}\StringTok{\textquotesingle{}obs\textquotesingle{}}\NormalTok{))}

\NormalTok{estadias\_totales }\SpecialCharTok{|\textgreater{}} \FunctionTok{head}\NormalTok{(}\DecValTok{15}\NormalTok{)}
\end{Highlighting}
\end{Shaded}

\begin{verbatim}
# A tibble: 15 x 8
   pais_de_residencia indice_tiempo turistas_no_residentes_BA
   <chr>              <date>                            <dbl>
 1 brasil             2014-01-01                       149582
 2 brasil             2014-04-01                       175018
 3 brasil             2014-07-01                       224259
 4 brasil             2014-10-01                       180361
 5 brasil             2015-01-01                       145922
 6 brasil             2015-04-01                       147540
 7 brasil             2015-07-01                       182629
 8 brasil             2015-10-01                       116157
 9 brasil             2016-01-01                        83492
10 brasil             2016-04-01                       127142
11 brasil             2016-07-01                       183799
12 brasil             2016-10-01                       145472
13 brasil             2017-01-01                       108533
14 brasil             2017-04-01                       146959
15 brasil             2017-07-01                       204501
# i 5 more variables: pernoctes_no_residentes_BA <dbl>,
#   estadia_media_no_residentes_BA <dbl>, turistas_no_residentes_CO <dbl>,
#   pernoctes_no_residentes_CO <dbl>, estadia_media_no_residentes_CO <dbl>
\end{verbatim}

\begin{Shaded}
\begin{Highlighting}[]
\FunctionTok{print}\NormalTok{(estadias\_totales }\SpecialCharTok{|\textgreater{}} \FunctionTok{select}\NormalTok{(pais\_de\_residencia) }\SpecialCharTok{|\textgreater{}} \FunctionTok{unique}\NormalTok{())}
\end{Highlighting}
\end{Shaded}

\begin{verbatim}
# A tibble: 3 x 1
  pais_de_residencia
  <chr>             
1 brasil            
2 chile             
3 nolimitrofe       
\end{verbatim}

Si observamos el \texttt{indice\_tiempo} las fechas parecieran repetirse
cada 3 meses, esto es porque los datos se actualizan trimestralmente, lo
cual podríamos haber anticipado si leíamos la página de donde
descargamos los datos.

\begin{Shaded}
\begin{Highlighting}[]
\NormalTok{estadias\_totales\_cuatrimestrales }\OtherTok{=}\NormalTok{ estadias\_totales }\SpecialCharTok{|\textgreater{}} 
  \FunctionTok{mutate}\NormalTok{(}
    \AttributeTok{anio=}\FunctionTok{year}\NormalTok{(indice\_tiempo),}
    \AttributeTok{trimestre=}\FunctionTok{case\_when}\NormalTok{(}
      \FunctionTok{month}\NormalTok{(indice\_tiempo) }\SpecialCharTok{==} \DecValTok{1}  \SpecialCharTok{\textasciitilde{}} \DecValTok{1}\NormalTok{,}
      \FunctionTok{month}\NormalTok{(indice\_tiempo) }\SpecialCharTok{==} \DecValTok{4}  \SpecialCharTok{\textasciitilde{}} \DecValTok{2}\NormalTok{,}
      \FunctionTok{month}\NormalTok{(indice\_tiempo) }\SpecialCharTok{==} \DecValTok{7}  \SpecialCharTok{\textasciitilde{}} \DecValTok{3}\NormalTok{,}
      \FunctionTok{month}\NormalTok{(indice\_tiempo) }\SpecialCharTok{==} \DecValTok{10} \SpecialCharTok{\textasciitilde{}} \DecValTok{4}\NormalTok{,}
      \AttributeTok{.default =} \ConstantTok{NA} \CommentTok{\# Opcional: para manejar meses que no coincidan}
\NormalTok{    )}
\NormalTok{  ) }\SpecialCharTok{|\textgreater{}} \FunctionTok{select}\NormalTok{(}\SpecialCharTok{!}\NormalTok{indice\_tiempo)}

\CommentTok{\# Versión tidy}
\NormalTok{estadias\_totales\_cuatrimestrales }\OtherTok{=}\NormalTok{ estadias\_totales }\SpecialCharTok{|\textgreater{}} 
  \FunctionTok{mutate}\NormalTok{(}
    \AttributeTok{anio=}\FunctionTok{year}\NormalTok{(indice\_tiempo),}
    \AttributeTok{trimestre=}\FunctionTok{quarter}\NormalTok{(indice\_tiempo)}
\NormalTok{  ) }\SpecialCharTok{|\textgreater{}} \FunctionTok{select}\NormalTok{(}\SpecialCharTok{!}\NormalTok{indice\_tiempo)}
\end{Highlighting}
\end{Shaded}

El objetivo de este bloque final es integrar herramientas de
programación, pensando en cómo sintetizar operaciones en un
\emph{script}.

\begin{itemize}
\item
  \textbf{Automatización mediante loops}: el uso de un bucle
  \texttt{for} permite escalar el análisis de forma eficiente. En lugar
  de repetir el código para cada mercado, generamos reportes
  estandarizados para Brasil, Chile y países No limítrofes con un solo
  bloque de ejecución.
\item
  \textbf{Jerarquía y factores}: la conversión de la columna
  \texttt{trimestre} a \textbf{factor} (\texttt{fct}) no es estética; es
  técnica. Garantiza que el orden cronológico (1 al 4) se preserve en
  cualquier análisis o gráfico posterior, evitando el ordenamiento
  alfabético por defecto.
\item
  \textbf{Persistencia de datos (.csv vs .rds)}: se guardan los datos en
  dos formatos distintos. El \texttt{.csv} permite la interoperabilidad
  con otras herramientas (como Excel), mientras que el \texttt{.rds} es
  muy útil en R porque es el único formato que guarda los
  \textbf{metadatos}, como los niveles de los factores que acabamos de
  configurar.
\end{itemize}

\begin{Shaded}
\begin{Highlighting}[]
\NormalTok{jerarquia }\OtherTok{=} \FunctionTok{c}\NormalTok{(}
  \StringTok{"trimestre"}\NormalTok{, }
  \StringTok{"porcentaje\_CO"}\NormalTok{,}
  \StringTok{"turistas\_no\_residentes\_BA"}\NormalTok{, }
  \StringTok{"turistas\_no\_residentes\_CO"}
\NormalTok{)}
\NormalTok{lista\_bases }\OtherTok{=} \FunctionTok{list}\NormalTok{(}
  \StringTok{"Brasil"} \OtherTok{=} \StringTok{\textquotesingle{}brasil\textquotesingle{}}\NormalTok{, }
  \StringTok{"Chile"} \OtherTok{=} \StringTok{\textquotesingle{}chile\textquotesingle{}}\NormalTok{, }
  \StringTok{"No limítrofes"}\OtherTok{=} \StringTok{\textquotesingle{}nolimitrofe\textquotesingle{}}
\NormalTok{ )}
\ControlFlowTok{for}\NormalTok{ (nombre }\ControlFlowTok{in} \FunctionTok{names}\NormalTok{(lista\_bases)) \{}
  \CommentTok{\# Nos quedamos con una tabla para cada categoría}
\NormalTok{  codigo\_pais\_residente }\OtherTok{=}\NormalTok{ lista\_bases[[nombre]]}
  
\NormalTok{  resumen }\OtherTok{=}\NormalTok{ estadias\_totales\_cuatrimestrales }\SpecialCharTok{|\textgreater{}} 
    \FunctionTok{filter}\NormalTok{(}
\NormalTok{      pais\_de\_residencia }\SpecialCharTok{==}\NormalTok{ codigo\_pais\_residente,}
\NormalTok{      anio}\SpecialCharTok{\textgreater{}=}\DecValTok{2015}\NormalTok{,}
      \DecValTok{2025}\SpecialCharTok{\textgreater{}=}\NormalTok{anio}
\NormalTok{    ) }\SpecialCharTok{|\textgreater{}} 
    \FunctionTok{select}\NormalTok{(}\SpecialCharTok{!}\FunctionTok{c}\NormalTok{(pais\_de\_residencia, anio)) }\SpecialCharTok{|\textgreater{}} 
    \FunctionTok{group\_by}\NormalTok{(}
\NormalTok{      trimestre}
\NormalTok{    ) }\SpecialCharTok{|\textgreater{}} 
    \FunctionTok{summarize}\NormalTok{(}
      \AttributeTok{turistas\_no\_residentes\_BA =} \FunctionTok{sum}\NormalTok{(}
\NormalTok{        turistas\_no\_residentes\_BA, }\AttributeTok{na.rm=}\ConstantTok{TRUE}
\NormalTok{      ),}
      \AttributeTok{turistas\_no\_residentes\_CO =} \FunctionTok{sum}\NormalTok{(}
\NormalTok{        turistas\_no\_residentes\_CO, }\AttributeTok{na.rm=}\ConstantTok{TRUE}
\NormalTok{      ),}
      \AttributeTok{.groups =} \StringTok{"drop"}
\NormalTok{    ) }\SpecialCharTok{|\textgreater{}} 
    \FunctionTok{mutate}\NormalTok{(}
      \AttributeTok{porcentaje\_CO =} \DecValTok{100}\SpecialCharTok{*}\NormalTok{turistas\_no\_residentes\_CO}\SpecialCharTok{/}\NormalTok{(turistas\_no\_residentes\_BA }\SpecialCharTok{+}\NormalTok{turistas\_no\_residentes\_CO)}
\NormalTok{    )}
  
  \CommentTok{\# Ordenamos las columnas jerárquicamente las columnas (esto NO es un factor)}
\NormalTok{  resumen }\OtherTok{=}\NormalTok{ resumen }\SpecialCharTok{|\textgreater{}} \FunctionTok{select}\NormalTok{(}\FunctionTok{all\_of}\NormalTok{(jerarquia))}
  
  \CommentTok{\# Convertimos \textquotesingle{}trimestre\textquotesingle{} en factor para asegurar orden cronológico (esto SÍ)}
\NormalTok{  resumen}\SpecialCharTok{$}\NormalTok{trimestre }\OtherTok{=} \FunctionTok{factor}\NormalTok{(resumen}\SpecialCharTok{$}\NormalTok{trimestre, }\AttributeTok{levels =} \DecValTok{1}\SpecialCharTok{:}\DecValTok{4}\NormalTok{)}
  
  \FunctionTok{print}\NormalTok{(}
    \FunctionTok{paste}\NormalTok{(}\StringTok{"Procesando:"}\NormalTok{, nombre)}
\NormalTok{  )}
  \FunctionTok{print}\NormalTok{(resumen)}
  
  \CommentTok{\# Definimos la ruta de salida}
  \FunctionTok{dir.create}\NormalTok{(}\StringTok{"../output"}\NormalTok{, }\AttributeTok{showWarnings =} \ConstantTok{FALSE}\NormalTok{) }\CommentTok{\# Crea la carpeta si no existe}
  
  \CommentTok{\# Creamos un nombre de archivo normalizado (ej: resumen\_brasil.csv)}
\NormalTok{  file\_name }\OtherTok{=} \FunctionTok{file.path}\NormalTok{(}\StringTok{"../output"}\NormalTok{, }\FunctionTok{paste0}\NormalTok{(}\StringTok{"resumen\_"}\NormalTok{, codigo\_pais\_residente))}
  
  \CommentTok{\# Exportación a CSV (Formato abierto)}
  \FunctionTok{write\_csv}\NormalTok{(resumen, }\FunctionTok{paste0}\NormalTok{(file\_name, }\StringTok{".csv"}\NormalTok{))}
  
  \CommentTok{\# Exportación a RDS (Formato nativo de R: preserva factores y tipos de datos)}
  \FunctionTok{saveRDS}\NormalTok{(resumen, }\FunctionTok{paste0}\NormalTok{(file\_name, }\StringTok{".rds"}\NormalTok{))}

\NormalTok{\}}
\end{Highlighting}
\end{Shaded}

\begin{verbatim}
[1] "Procesando: Brasil"
# A tibble: 4 x 4
  trimestre porcentaje_CO turistas_no_residentes_BA turistas_no_residentes_CO
  <fct>             <dbl>                     <dbl>                     <dbl>
1 1                  1.65                   1169674                     19646
2 2                  1.24                   1286842                     16186
3 3                  1.56                   1763087                     27974
4 4                  1.78                   1400661                     25412
[1] "Procesando: Chile"
# A tibble: 4 x 4
  trimestre porcentaje_CO turistas_no_residentes_BA turistas_no_residentes_CO
  <fct>             <dbl>                     <dbl>                     <dbl>
1 1                  4.83                    548407                     27843
2 2                  4.88                    496884                     25516
3 3                  5.49                    564384                     32780
4 4                  5.80                    600177                     36921
[1] "Procesando: No limítrofes"
# A tibble: 4 x 4
  trimestre porcentaje_CO turistas_no_residentes_BA turistas_no_residentes_CO
  <fct>             <dbl>                     <dbl>                     <dbl>
1 1                  2.76                   3995161                    113447
2 2                  3.02                   2425453                     75476
3 3                  3.36                   2347456                     81620
4 4                  2.88                   3507778                    103963
\end{verbatim}

Veamos cómo podemos cargar los datos y organizarlos juntos de forma
conveniente. La función \texttt{map\_df} de la librería \texttt{purrr}
(parte del \texttt{tidyverse}) automatiza la carga de múltiples
archivos:

\begin{itemize}
\item
  \textbf{\texttt{paises}}: es el vector con los nombres de los archivos
  a cargar (``brasil'', ``chile'', ``nolimitrofe'').
\item
  \textbf{\texttt{readRDS(...)}}: carga cada archivo en formato nativo
  de R desde la carpeta \texttt{../output}.
\item
  \textbf{\texttt{mutate(pais\ =\ .x)}}: agrega una columna nueva para
  identificar de qué archivo provienen los datos antes de unirlos.
\end{itemize}

Al terminar, \texttt{map\_df} une automáticamente todas las tablas en un
único dataframe consolidado.

\begin{Shaded}
\begin{Highlighting}[]
\CommentTok{\# Cargamos y unimos los datos de la carpeta output}
\NormalTok{paises }\OtherTok{=} \FunctionTok{c}\NormalTok{(}\StringTok{"brasil"}\NormalTok{, }\StringTok{"chile"}\NormalTok{, }\StringTok{"nolimitrofe"}\NormalTok{)}
\NormalTok{datos\_consolidados }\OtherTok{=} \FunctionTok{map\_df}\NormalTok{(}
\NormalTok{  paises, }
  \SpecialCharTok{\textasciitilde{}} \FunctionTok{readRDS}\NormalTok{(}\FunctionTok{file.path}\NormalTok{(}\StringTok{"../output"}\NormalTok{, }\FunctionTok{paste0}\NormalTok{(}\StringTok{"resumen\_"}\NormalTok{, .x, }\StringTok{".rds"}\NormalTok{))) }\SpecialCharTok{|\textgreater{}}  \FunctionTok{mutate}\NormalTok{(}\AttributeTok{pais =}\NormalTok{ .x) }\CommentTok{\# en este caso la tilde se usa para definir una formula que actua sobre .x (que, cuando itere paises, va a tomar el valor de cada país)}
\NormalTok{  )}
\end{Highlighting}
\end{Shaded}

Lo que sigue es necesario para hacer un gráfico, que no va a ser un
requisito del TP, es sólo para tener una idea de lo que anduvimos
haciendo.

\begin{Shaded}
\begin{Highlighting}[]
\CommentTok{\# Transformamos a formato LONG para separar por Provincia}
\NormalTok{datos\_histograma }\OtherTok{=}\NormalTok{ datos\_consolidados }\SpecialCharTok{|\textgreater{}} 
  \FunctionTok{pivot\_longer}\NormalTok{(}
    \AttributeTok{cols =} \FunctionTok{c}\NormalTok{(turistas\_no\_residentes\_BA, turistas\_no\_residentes\_CO), }\CommentTok{\# qué columnas colapsar}
    \AttributeTok{names\_to =} \StringTok{"provincia"}\NormalTok{, }\CommentTok{\# el nombre de la nueva columna que contendrá los antiguos nombres de las columnas }
    \AttributeTok{values\_to =} \StringTok{"total\_turistas"} \CommentTok{\# el nombre de la nueva columna que contendrá los valores numéricos que estaban en esas celdas}
\NormalTok{  )}
\end{Highlighting}
\end{Shaded}

Notar que al utilizar las funciones de la librería \texttt{forcats}
(\texttt{fct\_recode} y \texttt{fct\_reorder}) dentro del
\texttt{mutate}, transformamos esas columnas de texto plano
(\texttt{character}) a \textbf{factores} (\texttt{factor}), es decir,
que le decimos a R que conserve un orden jerárquico acorde a un
criterio.

\begin{Shaded}
\begin{Highlighting}[]
\FunctionTok{library}\NormalTok{(ggplot2)}
\FunctionTok{ggplot}\NormalTok{(datos\_histograma, }\FunctionTok{aes}\NormalTok{(}\AttributeTok{x =}\NormalTok{ trimestre, }\AttributeTok{y =}\NormalTok{ total\_turistas, }\AttributeTok{fill =}\NormalTok{ pais)) }\SpecialCharTok{+}
  \FunctionTok{geom\_bar}\NormalTok{(}\AttributeTok{stat =} \StringTok{"identity"}\NormalTok{, }\AttributeTok{position =} \StringTok{"dodge"}\NormalTok{) }\SpecialCharTok{+}
  \FunctionTok{facet\_wrap}\NormalTok{(}\SpecialCharTok{\textasciitilde{}}\NormalTok{ provincia, }\AttributeTok{scales =} \StringTok{"free\_y"}\NormalTok{) }\SpecialCharTok{+} \CommentTok{\# Crea un gráfico por provincia con su propia escala}
  \FunctionTok{theme\_minimal}\NormalTok{() }\SpecialCharTok{+}
  \FunctionTok{labs}\NormalTok{(}
    \AttributeTok{title =} \StringTok{"Distribución de Turistas por Orígen y Aeropuerto"}\NormalTok{,}
    \AttributeTok{subtitle =} \StringTok{"Comparativa trimestral entre Buenos Aires y Córdoba (2015{-}2025)"}\NormalTok{,}
    \AttributeTok{x =} \StringTok{"Trimestre"}\NormalTok{, }\CommentTok{\# el orden garantizado por que es factors}
    \AttributeTok{y =} \StringTok{"Cantidad Total de Turistas"}\NormalTok{,}
    \AttributeTok{fill =} \StringTok{"Origen"}
\NormalTok{  ) }\SpecialCharTok{+}
  \FunctionTok{scale\_fill\_brewer}\NormalTok{(}\AttributeTok{palette =} \StringTok{"Set2"}\NormalTok{)}
\end{Highlighting}
\end{Shaded}

\pandocbounded{\includegraphics[keepaspectratio]{data_wrangling_2_files/figure-pdf/unnamed-chunk-14-1.pdf}}

\subsection{4. Correr scripts en segundo
plano}\label{correr-scripts-en-segundo-plano}

Correr un script como \textbf{Background Job} en RStudio permite que R
ejecute el código en una sesión separada, dejándote la consola libre
para seguir analizando datos o probando funciones en la consola actual.
Para correr un script de esta forma hay que guardar el archivo e ir a
\texttt{Background\ jobs\ –\textgreater{}\ Start\ Background\ Jobround\ jobs\ –\textgreater{}\ Start\ Background\ Job}.
Seleccionamos el .R y el \textbf{directorio desde donde se corre.} De
esta forma, si estamos procesando un análisis exploratorio (como este),
pero en simultaneo queremos correr un archivo que puede tomar mucho
tiempo podemos ponerlo a correr en segundo plano.




\end{document}
